% =====================================================================
%  Appendix (TEXT-REVISER, revision 1): candidate data sources, the full
%  pre-registered analysis plan and the full ethics protocol, moved out of
%  Sections 2, 9 and 10 to shorten the main text. Every label is preserved.
% =====================================================================
\section{Candidate data sources}\label{app:data}

This appendix gives the source-by-source detail summarised in \cref{sec:data-incidents}.
For every source we record what can be checked in the provider's own documentation: the
timestamp the record carries, its temporal and geographic resolution, the anti-Jewish bias
category, and how the data can be obtained. Properties that we could not confirm from
primary documentation at the time of writing (November 2026) are marked as such and are not
used in any argument. \Cref{tab:data-streams} summarises the candidate streams.

\subsection{Incident streams}\label{app:data-incidents}

\paragraph{City-level open data: NYPD hate crimes.}
The New York City Police Department publishes a file of confirmed hate-crime
incidents on the NYC Open Data portal (dataset \texttt{bqiq-cu78})
\citep{nypdhate}, beginning in 2019. According to the portal's data
dictionary, \emph{Complaint Year Number} and \emph{Month Number} give the
year and month in which the incident occurred, and \emph{Record Create Date}
gives the date on which the report was filed. A \emph{Bias Motive
Description} field holds the NYPD's bias category, whose values include
\texttt{ANTI-JEWISH}. Geography is given by the complaint precinct, the
patrol borough and the county; each of the city's five counties is
coterminous with a borough. For point-process work the file thus has a
two-level time structure: the occurrence time is interval-censored to a
calendar month, while the report date is known to the day. We treat the
record-creation date as a report date and model the delay between occurrence
and record creation explicitly (\cref{sec:data-delays}), with the
occurrence month as a constraint. Whether day-level occurrence dates can be
obtained from the NYPD, for example under a data-use agreement, is an open
feasibility question for the pilot (\cref{sec:plan}); we do not assume it.
Changes to reporting systems are structural breaks for time-series use. The
New York State Comptroller reported that the state's Division of Criminal
Justice Services did not publish 2023 hate-crime data for New York City
because of the NYPD's transition in data-reporting systems as part of its
move to NIBRS \citep{osc2024hate}. \Citet{rosenbaum2025} analyses NYPD
administrative hate-crime records for 2019--2024 at monthly resolution.
Adjusting for the Jewish population of each borough, the study reports
that monthly anti-Jewish hate crimes were on average about twice as common
in the first year of the Israel--Hamas war as in the preceding years of the
study period (prevalence ratio 1.97, 95\% CI 1.64--2.35). Monthly
resolution suffices for such before--after contrasts but not for kernels
on the scale of days.

\paragraph{FBI NIBRS.}
The FBI's Uniform Crime Reporting Program retired its Summary Reporting
System on 1 January 2021 and moved to collection through the National
Incident-Based Reporting System (NIBRS) \citep{fbi2026nibrs}. For hate
crimes, however, the FBI reopened the 2021 collection to summary-format
submissions \citep{fbi2023supp}. According to the Congressional Research
Service, it has accepted both formats since 2022 \citep{crs2025hate}. Each
NIBRS incident carries the mandatory Data Element~3, \emph{Incident
Date/Hour}. The date is recorded as \texttt{YYYYMMDD} and the hour on a
24-hour clock. When the incident date is unknown, the report date is entered
instead and flagged with a report-date indicator
\citep{fbi2025nibrs,fbi2025nibrstech}. Anti-Jewish bias is code~21 of Data
Element~8A, \emph{Bias Motivation} \citep{fbi2025nibrs}. The geographic unit
is the reporting agency, identified by its ORI number. NIBRS submissions
thus supply occurrence dates, and hours where recorded, for individual
anti-Jewish incidents. For summary-format submissions we have not confirmed
comparable date fields, so the stream is restricted to NIBRS records. The
data are also slow. National statistics for data year 2024 were released on
5 August 2025 \citep{fbi2025release}, so for that year the lag from
occurrence to national availability was roughly 7 to 19 months, depending
on when in the year an incident occurred. Agency participation varies by
agency and year; the companion paper models it.

\paragraph{ADL Audit and H.E.A.T.\ Map.}
ADL describes its annual Audit as covering criminal and non-criminal
incidents of harassment, vandalism and assault, as reported to ADL by
victims, law enforcement, the media and partner organisations
\citep{adl2025audit}. An ADL stream therefore measures a different quantity
from a police stream, and the two reach the record through different funnels
(\cref{sec:data-thinning}). ADL states that since 7 October 2023 the Audit
has included certain expressions of opposition to Zionism and of support for
violent resistance against Israel or Zionists, while continuing to exclude
legitimate political protest \citep{adl2025audit}. For a point-process model
this matters for a structural reason, separate from any substantive
judgement. A change in recording scope at a date that is also a candidate
shock time $\tau_e$ is not identified separately from the shock term unless
the change is modelled explicitly. An ADL stream therefore needs a scope-break
indicator in $\mu_k(t)$, or records coded under a constant scope rule. ADL's
H.E.A.T.\ Map (Hate, Extremism, Antisemitism, Terrorism) publishes
incident-level listings that can be downloaded as a CSV file, with some
details removed to protect victims' privacy \citep{adlheatmap}. We have not
confirmed whether the listed date is the date of occurrence or of report to
ADL, nor how redaction affects dates. These checks are part of the pilot's
feasibility stage.

\paragraph{CST (United Kingdom).}
The Community Security Trust records antisemitic incidents in the United
Kingdom and publishes annual reports, most recently for 2025
\citep{cst2026incidents}. CST data would serve as an international comparison
stream. That report illustrates a further threat to time-series analysis.
CST states that since March 2024 a technical problem has disrupted the flow
of reports to it from the Metropolitan Police \citep{cst2026incidents}. In
the notation of \cref{sec:data-thinning}, an interruption in a referral feed
is a step change in the recording probability $p_k(t)$, which a model with
constant $p_k$ would read as a change in the incident process itself.
Incident-level records with dates would require a data-sharing agreement,
and a CST stream is therefore conditional.

\begin{table}[!htbp]
\centering
\caption{Candidate data streams for the pre-registered analysis. ``TBC'':
not confirmed from the provider's documentation at the time of writing; to
be confirmed before registration. Delays are from occurrence to
availability to analysts.}
\label{tab:data-streams}
\footnotesize
\setlength{\tabcolsep}{3pt}
\renewcommand{\arraystretch}{1.15}
\begin{tabularx}{\textwidth}{@{}>{\raggedright\arraybackslash}p{2.35cm}
  >{\raggedright\arraybackslash}X
  >{\raggedright\arraybackslash}p{2.2cm}
  >{\raggedright\arraybackslash}X
  >{\raggedright\arraybackslash}p{2.6cm}@{}}
\toprule
Stream & Time stamp (resolution) & Geography & Delay to availability & Access \\
\midrule
NYPD hate crimes \citep{nypdhate}
  & Occurrence: month; record creation: day. From 2019.
  & Precinct, patrol borough, county
  & Report delay bounded by occurrence month and record date; refresh schedule TBC
  & Public download (NYC Open Data) \\
FBI NIBRS \citep{fbi2025nibrs,fbi2026nibrs,crs2025hate}
  & Incident date and hour; report date substituted (flagged) if incident date unknown. NIBRS from 2021; hate-crime collection also accepts summary format since 2022.
  & Reporting agency (ORI)
  & Annual national release; data year 2024 released 5 Aug.\ 2025
  & Public (FBI Crime Data Explorer) \\
ADL Audit / H.E.A.T.\ Map \citep{adl2025audit,adlheatmap}
  & Incident date listed; occurrence vs.\ report date TBC. Scope change from 7 Oct.\ 2023.
  & Place listed; resolution TBC
  & TBC
  & Public map; CSV download with redactions \\
CST (UK) \citep{cst2026incidents}
  & TBC
  & TBC
  & Annual report; 2025 report published Feb.\ 2026
  & Incident-level records by agreement (TBC) \\
ACLED \citep{raleigh2010}
  & TBC
  & Event location
  & TBC
  & Provider terms (TBC) \\
GDELT \citep{leetaru2013}
  & TBC
  & Event location
  & TBC
  & Provider terms (TBC) \\
Platform data \citep{eu2022dsa}
  & Post time stamps (binned for analysis)
  & Varies
  & Depends on access route
  & DSA Art.~40 procedure; status TBC \\
Classifier volume \citep{kornsteinlab2026detection,kornsteinlab2026measurement}
  & Bin chosen by design (e.g.\ day)
  & None or coarse
  & Construction lag only
  & Lab-constructed from a defined corpus \\
\bottomrule
\end{tabularx}
\end{table}

\FloatBarrier
\subsection{Platform data access}\label{app:data-platform}

Article~40 of the EU Digital Services Act provides a procedure by which
vetted researchers may obtain data from very large online platforms and
very large online search engines for research that contributes to the detection,
identification and understanding of systemic risks in the Union. It also
contains a separate provision on access to publicly accessible data, and it
leaves technical conditions to delegated acts \citep{eu2022dsa}. Whether a
given project qualifies, and on what terms, is decided through that
procedure. We have not confirmed its operational status for the platforms
and periods of interest, and the plan in \cref{sec:plan} treats online
streams as optional, so that the offline analysis does not depend on them.
The terms of any research interface used will be recorded at registration.

\subsection{Prior empirical literature}\label{sec:data-lit}

\Citet{mohler2011} developed self-exciting point-process models of
crime; \citet{reinhart2018} reviews spatio-temporal versions and their
applications. Related applications include a spatiotemporal test of whether
gun violence is contagious \citep{loeffler2018} and self-exciting models of
civilian deaths in Iraq \citep{lewis2012}. For hate crime specifically,
\citet{king2013} study the temporal clustering of hate-motivated offending.
Two studies examine links between online content and offline hate crime:
\citet{muller2021} study social media and hate crime, and
\citet{williams2020} examine anti-Black and anti-Muslim social-media posts
as predictors of offline racially and religiously aggravated crime. For
anti-Jewish incidents, \citet{rosenbaum2025} compares monthly NYPD counts
before and after a salient external event (see above). We use this
literature to motivate the questions, not as prior estimates. None of these
studies estimates a multivariate branching matrix for antisemitic incidents,
and our framework does not take its parameter values from them. Separating
exogenous shocks from contagion is the identification problem of
\cref{sec:identification}.

\section{The pre-registered analysis plan in full}\label{app:plan}

This appendix sets out in full the analysis plan summarised in \cref{sec:plan}, which we
intend to register before any outcome data are examined. It fixes the candidate data, the
estimands, the procedure for listing external events, the model hierarchy and its tests,
the sensitivity analyses, the validation checks and the decision rules. The numerical
settings below (horizons, thresholds, grids, replicate numbers) are proposals. They will be
frozen at registration; any change made before registration will be documented, and any
change made after it will be reported as a deviation. The simulation study
(\cref{sec:simulation}) evaluates, on synthetic data with known ground truth, some of the
tools these rules rely on: EM estimation, surge attribution, the boundary test, the
thinning correction and a plug-in alert. It does not evaluate the block tests of
\cref{sec:plan-models}, the placebo calibration, the time-rescaling check, the Bayesian fit
or the selection and decision rules themselves; they are checked in the synthetic dry run
(V4). No part of the plan has yet been carried out on real data.

\subsection{Candidate pilot}\label{sec:plan-pilot}

\paragraph{Setting.}
The candidate pilot is New York City, using the NYPD hate-crime file
described in \cref{sec:data-incidents} \citep{nypdhate}. We restrict it
to records whose bias motive is \texttt{ANTI-JEWISH} and use the period from
2019 onwards. The offline streams $k$ are the five boroughs. Sparse streams
are handled by shrinkage on $\Amat$ (\cref{sec:estimation}), not by merging
streams on the basis of observed counts, which would let outcome data shape
the model. The offence category serves as a mark in a secondary analysis.
The final 12 months of data available at the freeze date form the hold-out
test period, and all earlier data form the training period.

\paragraph{Two designs, chosen before outcome data are seen.}
The NYPD file records the month of occurrence and the day of record creation
(\cref{sec:data-incidents}). The design is chosen on data availability
alone:
\begin{enumerate}[(A)]
  \item \emph{Design A} (conditional): if day-level occurrence dates can be
  obtained under a data-use agreement, the model is fitted to occurrence
  days, with within-day ties handled as in \cref{sec:data-needs}.
  \item \emph{Design B} (default): the model is fitted to record-creation
  dates, treated as a report-time process. The occurrence month enters as an
  interval constraint on each delay, and the delay distribution is modelled
  explicitly (\cref{sec:data-delays}). Contagion parameters then describe the
  report-time process, and sensitivity to the delay model is analysis~S4.
\end{enumerate}

\paragraph{Additional streams.}
Three further sources are conditional on access being confirmed during the
feasibility stage: (i) a classifier-based online volume series
(\cref{sec:data-online}); (ii) an ADL stream, if its export provides
day-level dates; and (iii) a CST stream as an international replication. A
national NIBRS analysis of anti-Jewish incidents is retrospective only,
because of the release lag (\cref{sec:data-incidents}). If a conditional
stream is unavailable it is dropped, and the primary analysis is unaffected.

\paragraph{Feasibility gates.}
Before registration we will: confirm the data dictionaries and archive the
data vintages (F1); choose between Designs A and B (F2); record the
licence terms of the event sources (F3); confirm or abandon the online
streams (F4); and obtain ethics approval (F5; \cref{sec:ethics}).

\subsection{Primary estimands}\label{sec:plan-estimands}

\begin{enumerate}[E1.]
  \item \emph{Branching matrix.} The matrix $\Amat=[\alpha_{kj}]$ for the
  recorded streams has three parts: the diagonal $\alpha_{kk}$
  (within-borough self-excitation); the off-diagonal offline entries
  (cross-borough excitation); and, when online streams are included, the
  online$\to$offline and offline$\to$online blocks. We report each entry
  with its interval, together with $\spr(\Amat)$ and
  $\Mmat=(\Imat-\Amat)^{-1}$. Estimates are reported on the recorded scale
  first. Thinning-corrected values (\cref{thm:thinning}) are reported
  alongside and labelled as model-based (\cref{sec:data-thinning}).
  \item \emph{Surge decomposition.} For each pre-registered event $e$ and
  stream $k$ we report the decomposition $\E[N_k(\tau_e,\tau_e+H]] =
  B_k + D_k + C_k$ of \cref{thm:attribution} and, where online streams are
  included, the online$\to$offline spillover share. The primary horizon is
  $H=14$ days, with $H\in\{7,28\}$ as secondary horizons. Results are
  summarised by event class. Event-level results appear in an appendix under
  neutral identifiers (\cref{sec:plan-events}).
  \item \emph{Alert performance.} We evaluate the rule of \cref{thm:alert}
  and \cref{alg:alert}, with window $\Delta_{\mathrm a}=7$ days. For each stream the
  surge threshold $c$ is the 95th percentile of the training-period
  distribution of 7-day counts. The cost ratio $c_{\mathrm{miss}}/c_{\mathrm{FA}}$
  is agreed with community-security partners before registration, with
  $\{2,5,10\}$ as the default grid. Metrics are the hit rate, false-alarm
  rate, realised cost and the Brier and log scores for the exceedance event.
  Two baselines serve as comparators: the background-only model M0 and a
  persistence rule that raises an alert when the previous window exceeded
  $c$.
\end{enumerate}
Secondary estimands are the kernel time scales, the class-level shock
loadings $\upsilon_k$ and the seasonal components of $\mu_k(t)$.

\subsection{Pre-specified external events}\label{sec:plan-events}

The event list is fixed by the following procedure. Every step is completed
before the analysis team loads any incident data for the study period.
\begin{enumerate}[1.]
  \item \emph{Classes.} Event classes are defined by type: (a) onset or
  major escalation of armed conflict; (b) terrorist or other mass-casualty
  attacks; (c) high-profile public statements by public officials or other
  prominent figures. The operational definitions are applied identically
  whatever the location, the parties, the perpetrators or the victims. We do
  not select events by their presumed relevance to antisemitism. Relevance
  is estimated through the loadings $\upsilon_k$, and events that are
  irrelevant to the outcome should receive loadings near zero. A secondary
  list restricts the pool to events whose coverage matches a keyword list
  that is fixed and deposited in advance.
  \item \emph{Salience threshold.} For each candidate event, a salience
  index is computed from the event datasets alone
  (\cref{sec:data-events}). It is the coverage attributed to the event over
  $[\tau_e,\tau_e+2\text{ days}]$, divided by the median daily coverage over
  the preceding 365 days. An event enters the list if its index exceeds a
  threshold $s^\star$. The threshold is set, using event data only, so that
  the list contains between 10 and 30 events over the study period. The
  size $S_e$ is the logarithm of the index.
  \item \emph{Blind coding.} Two coders independently review the candidates
  above the threshold. They merge duplicates, assign each event to a class
  and set $\tau_e$ to the time of first public report in the pilot city's
  time zone, to the day (or the hour under Design A). The coders have no
  access to incident data, which a separate data custodian holds until the
  freeze. Agreement is reported as Cohen's $\kappa$ for class and as the
  share of $\tau_e$ that agree to within one day. A third coder resolves
  disagreements. Events less than 48 hours apart are merged into a single
  shock at the earlier time.
  \item \emph{Freeze.} The list ($\tau_e$, class, $S_e$, source identifiers)
  is written to a file. The file's SHA-256 hash is posted with the
  registration and the file is deposited with the registry. Events after the
  freeze, which are needed for prospective evaluation, are coded under the
  same rules within seven days and appended with time-stamped hashes.
  \item \emph{Neutral identifiers.} Each event is reported by class and
  date. Descriptions follow the conventions of \cref{sec:ethics}.
\end{enumerate}

\subsection{Model hierarchy, block tests and selection rules}\label{sec:plan-models}

The models are nested:
\begin{enumerate}
  \item[M0] Background only: $\mu_k(t)=\exp(B(t)^\top\theta_k)$. The basis
  $B$ is a cubic B-spline in calendar time with knot spacing $h_b=60$ days, so that the
  number $K_b$ of basis functions grows with the length of the study period, plus
  day-of-week effects and indicators for a holiday calendar fixed in advance.
  \item[M1] M0 plus shocks $\sum_e\upsilon_k S_e\psi_k(t-\tau_e)$, with an
  exponential response kernel $\psi_k$; the loading $\upsilon_k$ and the
  kernel are specific to each stream and shared by all events of a class.
  \item[M2] M1 plus self-excitation (diagonal $\Amat$).
  \item[M3] M2 plus cross-stream offline excitation (the off-diagonal entries of the
  offline block of $\Amat$).
  \item[M4] M3 plus the online$\to$offline and offline$\to$online blocks, when online
  streams are available.
\end{enumerate}
The kernels are exponential, with one common decay rate within each block: the diagonal
(M2), the offline off-diagonal entries (M3), and each of the two online blocks (M4). We
estimate by maximum likelihood via EM (\cref{alg:em}), with parametric bootstrap intervals
from 500 replicates. A Bayesian fit with shrinkage on $\Amat$ (\cref{sec:est-bayes}) is
co-primary for M3 and M4.

\paragraph{Block tests.} Each step of the hierarchy adds a block of parameters that share
one rate: the loadings with their response rate (M0$\to$M1), or a block of entries of
$\Amat$ with its common decay (M1$\to$M2, M2$\to$M3, and for M4 first the online$\to$offline
block, then the offline$\to$online block). Under the null that every entry of the block is
zero, that rate is not identified, so the chi-bar-square reference
\eqref{eq:est-chibar} does not apply (\cref{rem:est-lrt}(i)). Each step is therefore decided
by the following registered test.
\begin{enumerate}[(i)]
  \item \emph{Statistic.} The sup-LR statistic \eqref{eq:est-suplr},
  $\mathrm{LR}_{\sup}=\max_g2\{\ell_T(\hat\vartheta_{j+1}(\beta_g))-\ell_T(\tilde\vartheta_j)\}$,
  where $\tilde\vartheta_j$ is the fit of the smaller model and
  $\hat\vartheta_{j+1}(\beta)$ the fit of the larger model with the new block's rate fixed
  at $\beta$ and every other parameter free.
  \item \emph{Grid.} $G=9$ rates whose half-lives $\log2/\beta_g$ are equally spaced on the
  log scale from $0.5$ to $21$ days ($0.5,\,0.8,\,1.3,\,2.0,\,3.2,\,5.2,\,8.3,\,13.2,\,21$
  days). The longest half-life is about a third of the knot spacing $h_b$, which keeps the
  cascade time scale apart from the background (\cref{sec:est-mle}).
  \item \emph{Reference.} The parametric bootstrap \eqref{eq:est-pboot} from the fit of the
  smaller model, with the observed calendar and $B=499$ replicates. In each replicate both
  models are refitted, the larger one at every grid value, starting from the null fit with
  small positive block entries (\cref{app:comp-tests}). The step is taken if
  $p_{\mathrm B}\le0.05$. The upcrossing approximation $p_{\mathrm D}$
  (\cref{app:est-davies}) is reported as a secondary quantity only.
  \item \emph{Order.} The block tests are applied in the fixed order of the hierarchy, each
  at level $0.05$ and only if the previous step was taken. This fixed-sequence procedure
  controls the family-wise error rate of the block decisions at $0.05$.
  \item \emph{Entries and the Holm family.} Within the blocks that were accepted, each entry
  $\alpha_{kj}$ is tested by the boundary test \eqref{eq:est-chibar}, whose decay is
  identified through the other entries of its block. If every other entry of the block is
  estimated as zero under the entry's null, that entry is tested with the sup-LR procedure
  of (i)--(iii) instead. All entry-level tests in the accepted blocks form one family, and
  Holm's procedure \citep{holm1979} controls its family-wise error rate at $0.05$. Decision
  rules D1 and D2 refer to these Holm-adjusted tests.
  \item \emph{Fallbacks.} The synthetic dry run (V4, \cref{sec:plan-validation}) checks the
  size and the cost of the block tests before any real data are loaded. If only the cost
  check fails, the score-based sup-statistic of \cref{app:est-davies} replaces the sup-LR
  statistic, with the same grid and $B=999$ replicates. If the size check fails, the block's
  rate is fixed in advance at the median of the grid (half-life $3.2$ days), and the block
  is tested by the likelihood-ratio statistic at that rate, with the parametric bootstrap
  reference of (iii). If that test also fails the size check, the step is decided by the
  out-of-sample log score alone, its excitation estimates are reported as exploratory, and
  no contagion claim under D1 or D2 is made for that block. The fallback used is reported
  for every step.
\end{enumerate}

\paragraph{Selection rules.} Moving from $\mathrm{M}j$ to $\mathrm{M}(j{+}1)$ requires both
of the following:
\begin{enumerate}[(a)]
  \item the block test of the step rejects at level $0.05$;
  \item the out-of-sample log score improves.
\end{enumerate}
The reported model is the largest model that meets both criteria. Estimates from every
model in the hierarchy are reported as well.

\subsection{Sensitivity analyses}\label{sec:plan-sens}

\begin{enumerate}[S1.]
  \item \emph{Background flexibility.} We vary the knot spacing
  $h_b\in\{30,60,120,365\}$ days and also fit a random-walk log-background.
  A misspecified, too rigid background produces spurious contagion
  (\cref{thm:spurious}), and second-order structure cannot settle the
  question (\cref{thm:equiv}). Self-excitation that vanishes under a more
  flexible background while goodness of fit remains acceptable is reported
  as not robustly distinguishable from background variation. An overly
  flexible background can absorb genuine contagion. In the simulation study a
  spline background removed the spurious contagion of slow regime switches but
  not that of fast unrecorded shocks (\cref{sec:sim-spurious}); how much
  genuine contagion it absorbs was not studied and is checked in the
  synthetic dry run (V4).
  \item \emph{Thinning.} We use (a) posterior draws of $p_k$ from the
  companion paper, mapped to streams as in \cref{sec:data-thinning}; (b) a
  grid $p_k\in\{0.1,0.25,0.5,0.75,1\}$; and (c) salience-dependent recording,
  $p_k(t)=p_k\{1+\varepsilon\sum_e\psi(t-\tau_e)\}$ with
  $\varepsilon\in\{0,0.25,0.5,1\}$ (capped at one), which violates the
  independence condition of \cref{thm:thinning}. In addition, (d) known
  recording breaks, such as reporting-system transitions, scope changes and
  interrupted referral feeds (\cref{sec:data-incidents}), enter $\mu_k(t)$
  as step indicators, and results are reported with and without the
  affected periods.
  \item \emph{Kernel family.} We compare exponential, gamma, power-law and
  nonparametric (histogram) kernels \citep{marsan2008}.
  \item \emph{Delay distribution.} Under Design B we fit the delay
  distribution in three versions: estimated from the month and
  record-creation bounds; with its scale shifted by $\pm50\%$; and with a
  heavy-tailed (log-normal) form. We compare the report-time analysis with
  an occurrence-month aggregate analysis. Nowcast corrections follow
  \cref{prop:nowcast}.
  \item \emph{Event list.} We halve and double $s^\star$, drop each event in
  turn, compare class-specific and pooled loadings, and vary the shock
  kernel $\psi$.
  \item \emph{Spatial resolution.} We compare boroughs with patrol boroughs.
  \item \emph{Online measurement error.} We compare uncorrected and corrected
  series (\cref{eq:data-classerr}), with period-specific error rates and
  propagated uncertainty.
  \item \emph{Ties.} We compare at least 20 within-day randomisations with
  the discrete-time INAR fit (\cref{eq:inar}).
\end{enumerate}

\subsection{Validation}\label{sec:plan-validation}

\begin{enumerate}[V1.]
  \item \emph{Goodness of fit.} Following \citet{brown2002}, we apply the
  time-rescaling theorem (\cref{prop:rescaling}) stream by stream: a
  Kolmogorov--Smirnov test on rescaled inter-event times, with times
  randomised within the recorded resolution. Because the parameters are
  estimated, the reference distribution is obtained by parametric
  bootstrap.
  \item \emph{Out-of-sample performance.} We compute log scores for daily
  counts \citep{gneiting2007} and randomised PIT histograms for calibration
  \citep{dawid1984,gneiting2007b}, using rolling-origin one-step and
  $\Delta_{\mathrm a}$-step forecasts over the test period.
  \item \emph{Placebo shocks.} Placebo calibration is done for one event class at a time
  (\cref{sec:est-placebo}). For each class we draw 1{,}000 placebo calendars with the same
  number of events as the real class, and permute the real sizes over them. Their dates are
  drawn uniformly from \emph{all} days of the training period, holidays included, subject
  only to the 48-hour merging rule of \cref{sec:plan-events} within each calendar, and
  without reference to the real dates of that class. Excluding holidays, or any other set of
  days on which a real event could fall, would break the exchangeability on which the rank
  test rests. The base model is the selected model with that class removed. The real class
  and each placebo calendar are added to the base model in turn, as a shock class, and the
  same statistics are computed for all of them. The rejection frequency of the placebo
  loading tests at nominal level $0.05$ is the empirical false-positive rate of the shock
  test. The rank of each real-event direct component $D_k$ among its placebo counterparts is
  reported. The rank $p$-value is exact under the randomisation hypothesis of
  \cref{lem:est-placebo}, that the real dates are exchangeable with draws from this design;
  otherwise it is reported as a calibration, not as a test. If the empirical
  false-positive rate exceeds $0.10$, inference on shocks uses the placebo distribution in
  place of the asymptotic reference.
  \item \emph{Synthetic dry run.} Before any real data are loaded, the full pipeline is run
  on synthetic data generated as in \cref{sec:simulation}, at the pilot's number of
  streams, length of study period and calendar size. It uses $R=200$ datasets from each of
  two configurations: a null configuration without excitation (M1), and an alternative
  with self-excitation $0.3$ and cross-stream entries $0.05$, both with a spline
  background. It checks (a) the size of each block test of \cref{sec:plan-models} at a
  configuration that satisfies its null (the null configuration for the excitation blocks,
  and the same configuration without shocks for M0$\to$M1), which passes if the rejection
  rate is at most $0.075$ (Monte Carlo standard error about $0.015$); (b) its cost, which passes if one block test takes at
  most seven days on the project's computing allocation; (c) the coverage of the intervals
  for $\alpha_{kj}$ and for the surge decomposition at the alternative; and (d) the
  selection and decision rules, including whether a spline background with spacing $h_b$
  absorbs genuine contagion. Failures of (a) or (b) trigger the fallbacks of
  \cref{sec:plan-models}.
\end{enumerate}

\subsection{Decision rules}\label{sec:plan-decisions}

\begin{enumerate}[D1.]
  \item \emph{Self-excitation} in stream $k$ is claimed only if: the test of
  $\alpha_{kk}=0$ rejects after the Holm correction; the 95\% interval for
  $\alpha_{kk}$ excludes zero under the most flexible background in S1;
  goodness of fit (V1) is not rejected at level $0.01$; and the placebo
  false-positive rate (V3) is at most $0.10$. Otherwise we report that
  self-excitation is not distinguishable from background variation at the
  available resolution.
  \item \emph{Cross-stream and online$\to$offline excitation} are claimed
  under the same criteria applied to $\alpha_{kj}$. They are described as
  Granger-type predictive dependence (\cref{prop:granger}) and never as
  causation.
  \item \emph{Surge decompositions} are reported for every pre-registered
  event, whatever its size or significance. A surge is described as
  predominantly due to contagion only if the 95\% interval for
  $C_k/(D_k+C_k)$ lies above $0.5$ under the primary model and under the
  S1 and S2 variants.
  \item \emph{Thinning-corrected} branching ratios are never reported
  without the $p_k$ used and the S2 range.
  \item \emph{Alerts} are recommended for operational use with partners only
  if, in the test period, the rule's realised cost is lower than that of
  both baselines, with a bootstrap 95\% interval for the difference that
  excludes zero, and the PIT shows no marked miscalibration. Otherwise the
  alert analysis is reported as a research result only.
  \item \emph{Null results and deviations} are reported in full.
\end{enumerate}

\subsection{Replication package}\label{sec:plan-replication}

The package will contain:
\begin{enumerate}[(i)]
  \item the estimation and simulation code with a pinned environment;
  \item the frozen event list and its hash;
  \item acquisition scripts with dataset identifiers, vintages, download
  dates and SHA-256 checksums;
  \item derived counts by stream, at the resolutions and subject to the
  small-cell rules of \cref{sec:ethics-irb};
  \item for restricted streams (platform or partner data), the code together
  with a synthetic dataset of the same structure;
  \item all results as machine-readable files, with scripts that regenerate
  every table and figure; and
  \item the deviations log.
\end{enumerate}
Raw records are redistributed only where the provider's terms permit.

\subsection{Timeline}\label{sec:plan-timeline}

\Cref{tab:plan-timeline} gives the timeline in months from the start of the
feasibility stage.

\begin{table}[!htbp]
\centering
\caption{Timeline of the pre-registered analysis, in months from the start
of the feasibility stage. Each gate must be passed before the next phase
begins.}
\label{tab:plan-timeline}
\footnotesize
\renewcommand{\arraystretch}{1.15}
\begin{tabularx}{\textwidth}{@{}l l >{\raggedright\arraybackslash}X >{\raggedright\arraybackslash}p{4.2cm}@{}}
\toprule
Phase & Months & Activities & Gate / output \\
\midrule
P0 Feasibility & 0--2 & Confirm data dictionaries and access terms (F1--F4); choose Design A or B; submit for ethics review & Ethics approval; design fixed \\
P1 Events and registration & 2--3 & Blind coding; freeze event list; register plan & Time-stamped registration; hash deposited \\
P2 Dry run & 3 & Full pipeline on synthetic data (V4) & Coverage and selection rules checked \\
P3 Training fit & 4--6 & Fit M0--M4; apply selection rules & Selected model \\
P4 Test evaluation & 7 & V1--V3; back-test of alerts & Out-of-sample report \\
P5 Sensitivity & 7--9 & S1--S8 & Robustness tables \\
P6 Reporting & 9--11 & Paper; replication package; partner debrief under \cref{sec:ethics} & Public release \\
P7 Prospective & 12--24 & Prospective alert evaluation under embargo & Embargoed evaluation report \\
\bottomrule
\end{tabularx}
\end{table}

\FloatBarrier
\section{The ethics protocol in full}\label{app:ethics}

This appendix gives in full the protocol summarised in \cref{sec:ethics}. It accompanies
the plan of \cref{app:plan} and complements the companion paper's protocol for victim
privacy and record linkage \citep{kornsteinlab2026darkfigure}.

\subsection{Risks of forecasting}\label{sec:ethics-risks}

\paragraph{Stigmatisation.}
A map of elevated predicted intensity can be read as a statement about the
people who live in an area. That reading applies to neighbourhoods where
Jewish residents are targeted, which can come to be seen as ``dangerous''.
It applies equally to any community that observers might wrongly treat as a
source of incidents. The model contains no information about perpetrators,
and its spatial units are recording areas, not communities. Alerts and published maps are therefore
at the borough level or coarser, and never street- or venue-level heat maps.
Finer units used in sensitivity analysis S6 are reported only as summary
statistics. No output ranks neighbourhoods or communities.

\paragraph{Self-fulfilling and self-defeating alerts.}
An alert can change what it predicts. Publicity can raise reporting, which
inflates recorded counts (the time-varying $p_k$ of analysis S2). Protective
deployment can prevent incidents, which makes a correct alert look like a
false alarm. Either effect contaminates prospective evaluation. Alerts
during the prospective period (phase P7) are therefore shared only under the
protocol below, and evaluation reports state which alerts were acted on.

\paragraph{Information hazards.}
Forecasts that name ``hot'' dates or locations could help someone planning
an attack, and could encourage imitation. We therefore publish methods,
code and retrospective results, but no live or near-live forecasts and no
forecast tied to a specific site. Retrospective results are released only
after an embargo (\cref{sec:ethics-disclosure}).

\subsection{Responsible disclosure of alerts}\label{sec:ethics-disclosure}

Alerts produced during the prospective evaluation are handled as follows.
\begin{enumerate}[1.]
  \item \emph{Recipients.} Alerts go only to vetted community-security
  partners, and to a police liaison where the partner agreement provides for
  this. Each partner signs an agreement that limits use to protective
  purposes, such as security staffing and victim support, and excludes
  enforcement targeting of any community.
  \item \emph{Content.} An alert gives the borough or coarser area, the
  window $\Delta_{\mathrm a}$, the predicted probability of exceeding the threshold
  with its uncertainty, and a fixed statement of the model's limitations.
  It contains no individual-level information and no attribution to
  persons, groups or communities.
  \item \emph{Human review.} A named member of the research team reviews
  every alert before release, and may withhold an alert that is plainly
  driven by a data artefact.
  \item \emph{Channel and logging.} Alerts are sent over an encrypted
  channel. Each alert, the decision to release or withhold it, and the
  partner's acknowledgement are logged for audit.
  \item \emph{Embargo.} Alerts are not published. Aggregate evaluation
  results (hit and false-alarm rates by quarter and borough) are published
  no earlier than 90 days after the end of the period they cover.
  \item \emph{Suspension.} Alerts are suspended if rolling calibration
  deteriorates beyond a pre-specified tolerance, or if a partner reports
  misuse.
\end{enumerate}

\subsection{Language and attribution}\label{sec:ethics-language}

The model attributes \emph{timing}, not blame. A positive shock loading
$\upsilon_k$, or a large direct component $D_k$ after event $e$, means that
recorded incidents in stream $k$ rose after $\tau_e$ by more than the
background and the contagion terms predict. It does not mean that the
parties to the event, the communities associated with them, or the people
who discussed it caused any incident. Incidents are committed by
individual perpetrators, about whom the model contains no information. We
therefore:
\begin{enumerate}[(i)]
  \item describe events as \emph{salient external events} of a stated class
  and date, and do not characterise them politically;
  \item never describe a community, nationality or religious group as a
  source, driver or cause of incidents;
  \item phrase every result in terms of recorded incidents and statistical
  association; this includes the online$\to$offline coefficients, which
  are Granger-type (\cref{prop:granger}); and
  \item report event-level decompositions in an appendix under neutral
  identifiers, with class-level summaries in the main text.
\end{enumerate}

\subsection{Political neutrality}\label{sec:ethics-neutrality}

The paper takes no position on any political conflict, government, movement
or public figure. Neutrality is built into the procedure, not just asserted:
event classes are defined by type and apply to all parties alike
(\cref{sec:plan-events}); the event list is frozen before outcome data are
seen, so events cannot be selected after the fact; every pre-registered
event is reported whatever the result (rule D3); and the event-coding rules
and frozen list are public. Before registration, the event-coding protocol
will be reviewed by at least two external reviewers with differing
perspectives on the events likely to be coded. Their comments and our
responses will be deposited with the registration.

\subsection{Online data}\label{sec:ethics-online}

Online data are used only under the terms of the access route through which
they are obtained, including any conditions on data protection and security
attached to access under Article~40 of the Digital Services Act
\citep{eu2022dsa}. The analysis uses aggregate counts per time bin. We do
not build account-level scores or profiles, do not link online accounts to
offline incidents, victims or suspects, and do not publish post text,
handles or quotations that would allow re-identification. Annotation
samples for estimating classifier error (\cref{eq:data-classerr}) are held
under access control and deleted when the retention period agreed with the
data provider and the ethics board ends. Classifier outputs are used only as
aggregate series. They are never used to flag individuals.

\subsection{Ethics review and data security}\label{sec:ethics-irb}

The protocol, including this section, will be submitted to an institutional
review board or equivalent research ethics committee before any data are
accessed (gate F5). This applies even to publicly downloadable data, because
the plan combines several streams and may involve partner and platform
data. Restricted data are kept on encrypted storage with role-based access,
and use is logged. The custodian who holds the incident data during event
coding (\cref{sec:plan-events}) is separate from the coders. Outputs follow
a minimum-cell rule: counts below 10 are not shown below the borough level,
and any table that could reveal a single incident's date and location is
suppressed or coarsened. Derived counts are released at a resolution no
finer than the source itself publishes; for restricted streams, they are
released weekly or more coarsely. Data are retained no longer than the agreements
allow, and the replication package contains only data whose release is
permitted (\cref{sec:plan-replication}).
